% Author's intro paragraph as edited October 26, 2026, 19:39 ET, before the slide-spine version replaced it in the mock.

The decline in fertility stands out as one of the most striking phenomena taking place in today's world.   
Fertility in the United States  fallen steadily globally, and the
total fertility rate in the United States stands today at about 1.6 children per
woman.\footnote{The total fertility rate peaked at 2.12 in 2007 and was 1.60
in 2024 \citep{cboDemographicOutlook2026}. The causes of the decline are
debated. \citet{kearneyPuzzleFallingUS2022} find no economic or policy
factor that accounts for it. \citet{kearneyLevineWhyFertilityLow2025} argue
for ``shifting priorities'' across cohorts, a change in norms around work,
parenting, gender roles and leisure, shaped in part by new technologies such
as the smartphone and social media.
\citet{doepkeEconomicsFertilityNew2023} emphasize how easily parents can
combine careers and family.} Low fertility raises concerns well beyond
family life: fewer workers support each retiree, public pension and health
systems come under strain, and population growth comes to rest on
immigration alone.\footnote{The Congressional Budget Office projects that
deaths will exceed births from 2030 \citep{cboDemographicOutlook2026}. The
Social Security Trustees lowered their long-run assumption for the total
fertility rate from 1.90 to 1.75 children per woman \citep{ssaTrustees2026}.}
In the public debate, housing costs are often at the center of the
conversation, from surveys of prospective parents to recent proposals for
down-payment assistance.\footnote{In a fourteen-country survey, housing is
the third most cited reason for having fewer children than desired, named by
19 percent of respondents \citep{unfpaSWP2025}. In 2024, the Harris campaign
proposed up to \$25,000 in down-payment assistance for first-time buyers
\citep{harrisCampaign2024}.} The reasons are intuitive. Children need space,
and new parents want extra rooms. Young households, who are the ones having
children, have little savings, so the affordability of that space matters
most to them. This paper asks how the cost of family space shapes fertility,
and what housing policy can do about it.
